% ============================================================
% Section 4.4: The Asynchronous Gauss-Markov Theorem
%
% This section establishes that the two-step KKT estimator is
% the BLUE within the class of all unbiased cross-product
% estimators of sigma_12 from asynchronous tick data.
%
% Three claims are established:
%   Claim 1 (Lemma 4.1): characterisation of the estimator class
%   Claim 2 (proved in Theorem 2): KKT minimises variance
%   Claim 3 (Theorem 4): equality with CRLB at rho=0
% ============================================================

\subsection{The Asynchronous Gauss-Markov Theorem}
\label{sec:gaussmarkov}

We now establish the optimality of the two-step estimator in its
strongest form.  The classical Gauss-Markov theorem identifies the
best linear unbiased estimator (BLUE) within the class of linear
estimators.  We prove an analogue for the asynchronous covariance
estimation problem, identifying the two-step estimator as the BLUE
within the natural class of unbiased estimators of~$\sigma_{12}$
from asynchronous tick data.

\paragraph{The natural estimator class.}
The observations are the return vectors
$\mathbf{R} = (R_1,\dots,R_m)^\top$ and
$\mathbf{S} = (S_1,\dots,S_n)^\top$.
Since $\EE R_i = \EE S_j = 0$, any estimator that is linear in the
individual returns has zero expectation and cannot be unbiased for
$\sigma_{12} \neq 0$.  The natural estimator class therefore
consists of \emph{quadratic} functions of the observations --
specifically, linear combinations of the cross-products $R_iS_j$.

\begin{lemma}[Characterisation of unbiased estimators]
\label{lem:unbiased_class}
Let $\hat\sigma = \sum_{i,j} c_{ij} R_i S_j$ be any estimator of
$\sigma_{12}$ that is a linear combination of cross-products.
Then $\hat\sigma$ is unbiased for $\sigma_{12}$ if and only if
\begin{equation}\label{eq:unbiased_char}
    \sum_{\{(i,j):\,\tau_{ij}>0\}} c_{ij}\,\tau_{ij} \;=\; 1.
\end{equation}
In particular, the only cross-products $R_iS_j$ that contribute to
an unbiased estimator are those for which $\tau_{ij}>0$, i.e.\
those corresponding to \emph{overlapping} return intervals.
\end{lemma}

\begin{proof}
Under the model~\eqref{eq:bivariateBM} with constant parameters,
$\EE[R_iS_j] = \sigma_{12}\,\tau_{ij}$ for all $i,j$.  Therefore
$\EE\hat\sigma = \sigma_{12}\sum_{ij}c_{ij}\,\tau_{ij}$, which
equals $\sigma_{12}$ if and only if~\eqref{eq:unbiased_char} holds.
Since $\tau_{ij}=0$ for non-overlapping pairs, those terms contribute
nothing to either the expectation or the constraint, and may be set
to zero without loss.
\end{proof}

\begin{remark}
Lemma~\ref{lem:unbiased_class} can also be derived from the
structure of the Fisher score.  The full joint distribution of
$(\mathbf{R},\mathbf{S})$ is multivariate Gaussian with covariance
matrix $\Sigma(\sigma_{12})$ in which $\sigma_{12}$ appears only
in the off-diagonal blocks: $[\partial\Sigma/\partial\sigma_{12}]_{i,m+j} = \tau_{ij}$
and zero elsewhere.  The score for $\sigma_{12}$ is therefore
$\frac{1}{2}\mathbf{y}^\top\Sigma^{-1}
(\partial\Sigma/\partial\sigma_{12})\Sigma^{-1}\mathbf{y}
- \frac{1}{2}\operatorname{tr}(\Sigma^{-1}\partial\Sigma/\partial\sigma_{12})$,
which is a quadratic form in $\mathbf{y}$ involving only the
cross-product terms.  Any unbiased estimator must satisfy the same
structural constraint.
\end{remark}

Lemma~\ref{lem:unbiased_class} shows that the class of unbiased
estimators considered in this paper -- the weighted ZHY estimators
$\hat\sigma_{12}(\mathbf{w})$ -- is \emph{complete}: it exhausts
all unbiased cross-product estimators of $\sigma_{12}$.
The two-step estimator therefore minimises variance over the full
natural estimator class, not merely over a restricted subclass.

\begin{thm}[Asynchronous Gauss-Markov Theorem]
\label{thm:gaussmarkov}
Under the assumptions of Section~\ref{sec:cross} with constant
parameters $\sigma_1$, $\sigma_2$, $\sigma_{12}$, the two-step
estimator $\hat\sigma_{12}^{(2)}$ of Definition~\ref{def:twostep}
is the \emph{best linear unbiased estimator} (BLUE) of $\sigma_{12}$
within the class of all unbiased cross-product estimators
$\sum_{ij}c_{ij}R_iS_j$ satisfying~\eqref{eq:unbiased_char}.
That is,
\begin{equation}\label{eq:gaussmarkov}
    \VV\!\left(\hat\sigma_{12}^{(2)}\,\big|\,\tau,\hat\kappa\right)
    \;\leq\;
    \VV\!\left(\sum_{ij}c_{ij}R_iS_j\,\Big|\,\tau\right)
\end{equation}
for every coefficient vector $\{c_{ij}\}$
satisfying~\eqref{eq:unbiased_char}.

Furthermore, at $\rho = 0$, the two-step estimator achieves the
Cram\'er-Rao lower bound:
\begin{equation}\label{eq:cramer_rao_zero}
    \VV\!\left(\hat\sigma_{12}^{(2)}\,\big|\,\tau,\,\kappa=0\right)
    \;=\; \frac{1}{I(\sigma_{12}|\tau)}\,,
\end{equation}
where $I(\sigma_{12}|\tau)$ is the Fisher information for
$\sigma_{12}$ given the observed timestamps.
\end{thm}

\begin{proof}
The first claim~\eqref{eq:gaussmarkov} follows from
Lemma~\ref{lem:unbiased_class} and Theorem~\ref{thm:twostep}.
By Lemma~\ref{lem:unbiased_class}, every unbiased cross-product
estimator of $\sigma_{12}$ is of the form $\hat\sigma_{12}(\mathbf{c})
= \sum_{ij}c_{ij}R_iS_j$ satisfying~\eqref{eq:unbiased_char}.
Re-writing in the normalised form $\hat\sigma_{12}(\mathbf{w})$
of~\eqref{eq:genZHYest}, the class of such estimators coincides
with the weighted ZHY class.  By Theorem~\ref{thm:twostep}, the
two-step estimator minimises variance over this class, which
establishes~\eqref{eq:gaussmarkov}.

For the second claim~\eqref{eq:cramer_rao_zero}, set $\sigma_{12}=0$
(equivalently $\kappa=0$).  At $\kappa=0$ the KKT
system~\eqref{eq:kkt} reduces to $w_{ij}^*\tau_i^X\tau_j^Y
= \mu\,\tau_{ij}$, giving
$w_{ij}^* = \tau_{ij}/(\tau_i^X\tau_j^Y)$.  The corresponding
variance is $\sigma_1^2\sigma_2^2 A_1^*$, where
\begin{equation*}
    A_1^* \;=\; \frac{\displaystyle\sum_{ij}
                      (w_{ij}^*)^2\tau_i^X\tau_j^Y}
                     {\displaystyle\left(\sum_{ij}\tau_{ij}w_{ij}^*\right)^2}
             \;=\; \frac{1}{\displaystyle\sum_{ij}\tau_{ij}^2/(\tau_i^X\tau_j^Y)}\,.
\end{equation*}
When $\sigma_{12}=0$, the observations $\mathbf{R}$ and $\mathbf{S}$
are independent.  The joint covariance matrix is block-diagonal:
$\Sigma = \operatorname{diag}(\sigma_1^2\tau_i^X,\,\sigma_2^2\tau_j^Y)$.
The Fisher information for $\sigma_{12}$ is
\begin{align*}
    I(\sigma_{12}|\tau)\big|_{\sigma_{12}=0}
    &= \tfrac{1}{2}\operatorname{tr}\!\left[
        \Sigma^{-1}\frac{\partial\Sigma}{\partial\sigma_{12}}
        \Sigma^{-1}\frac{\partial\Sigma}{\partial\sigma_{12}}
       \right]_{\sigma_{12}=0} \\
    &= \tfrac{1}{2}\sum_{i,j}
       \frac{2\,\tau_{ij}^2}{\sigma_1^2\tau_i^X \cdot \sigma_2^2\tau_j^Y}
     \;=\; \frac{\displaystyle\sum_{ij}\tau_{ij}^2/(\tau_i^X\tau_j^Y)}
                {\sigma_1^2\sigma_2^2}\,,
\end{align*}
where we used $[\Sigma^{-1}]_{ii} = 1/(\sigma_1^2\tau_i^X)$,
$[\Sigma^{-1}]_{m+j,m+j} = 1/(\sigma_2^2\tau_j^Y)$,
and off-diagonal blocks zero.
Therefore $1/I(\sigma_{12}|\tau)|_{\sigma_{12}=0}
= \sigma_1^2\sigma_2^2/\sum_{ij}\tau_{ij}^2/(\tau_i^X\tau_j^Y)
= \sigma_1^2\sigma_2^2 A_1^*$,
which is precisely the variance of the two-step estimator.
This establishes~\eqref{eq:cramer_rao_zero}.
\end{proof}

\begin{remark}[Information-theoretic interpretation]
Theorem~\ref{thm:gaussmarkov} has a natural information-theoretic
interpretation.  At $\rho=0$, the cross-products $\{R_iS_j\}$ are
\emph{sufficient statistics} for $\sigma_{12}$: because $\mathbf{R}$
and $\mathbf{S}$ are independent, the joint likelihood factorises
and all information about $\sigma_{12}$ is carried by the
cross-covariance structure.  The ZHY estimator therefore exhausts
the available information and achieves the CRLB.

At $\rho\neq 0$, the marginal variances of $R_i$ and $S_j$ also
carry information about $\sigma_{12}$ (through their joint
distribution).  This information is inaccessible to any linear
cross-product estimator, creating the gap between the ZHY class
and the CRLB visible in Table~\ref{tab:subopt}.  Closing this gap
requires a nonlinear estimator -- the maximum likelihood estimator --
which involves inverting the full covariance matrix $\Sigma(\sigma_{12})$.
This direction is discussed in Section~\ref{sec:discussion} as an
avenue for future work.
\end{remark}

\begin{remark}[Relationship to the classical Gauss-Markov theorem]
The classical Gauss-Markov theorem applies to the linear model
$\mathbf{y} = X\boldsymbol\beta + \boldsymbol\varepsilon$, where
the BLUE of $\boldsymbol\beta$ is the OLS estimator (under
homoscedasticity) or the GLS estimator (under heteroscedasticity).
Theorem~\ref{thm:gaussmarkov} is the analogue for a
\emph{second-order} parameter $\sigma_{12}$ in a multivariate
Gaussian model with non-standard covariance structure induced by the
asynchronous observation scheme.  The role of OLS is played by the
basic ZHY estimator; the role of GLS is played by the two-step
KKT estimator.  The Gauss-Markov property is established by the
same convexity argument as in the classical case, applied to the
quadratic form $V(\mathbf{w}) = \sigma_1^2\sigma_2^2 A_1(\mathbf{w})
+ \sigma_{12}^2 A_2(\mathbf{w})$.
\end{remark}

\begin{remark}[Numerical verification]
The equality~\eqref{eq:cramer_rao_zero} was verified numerically
across 100 randomly generated timestamp structures with dimensions
$m \in \{2,3,4\}$, $n \in \{2,3,4\}$.  In every case the relative
error $|\VV(\hat\sigma_{12}^{(2)}|\tau) - 1/I(\sigma_{12}|\tau)|
\,/\,I(\sigma_{12}|\tau)^{-1}$ was below $10^{-14}$, consistent
with floating-point precision.  This provides strong numerical
confirmation of the analytic result.
\end{remark}
